\documentclass[aps,prb,twocolumn,floatfix,nofootinbib,superscriptaddress]{revtex4-2}

\usepackage{amsmath,amssymb,amsthm,mathtools,bm}
\usepackage{booktabs}
\usepackage[colorlinks=true,linkcolor=blue,citecolor=blue,urlcolor=blue]{hyperref}
\usepackage{microtype}

\newcommand{\tr}{\operatorname{tr}}
\newcommand{\ii}{\mathrm{i}}
\newcommand{\code}[1]{\texttt{#1}}

\hypersetup{
  pdftitle={Protocol for Frozen-Record Boundary-Twist Charge Redistribution in an Adaptive Chern Circuit},
  pdfauthor={Asadullah Bhuiyan}
}

\begin{document}

\title{\texorpdfstring{Protocol for Frozen-Record Boundary-Twist Charge
Redistribution\\in an Adaptive Chern Circuit}{Protocol for Frozen-Record
Boundary-Twist Charge Redistribution in an Adaptive Chern Circuit}}
\author{Asadullah Bhuiyan}
\affiliation{Class-A fermionic adaptive-circuit project}
\date{September 3, 2026}

\begin{abstract}
We specify a compact CPU pilot for diagnosing boundary-twist charge
redistribution in the class-A adaptive free-fermion circuit.  A zero-twist
Born record is generated once and then imposed in independent static replays
at a grid of twisted boundary conditions.  Separate records are required for
soft explicit interfaces and hard support-terminated walls because their
active measurement sets differ.  The primary geometry is a
$16\times20$ cylinder evolved for $2N_y=40$ perfect-correction cycles with a
pure half-filled initial state, random serial measurement order,
$n_{\mathrm{shell}}=1$, and complex128 arithmetic.  Both increasing and
decreasing sweeps use 17 points with the signed origin convention
$\phi_j^{(\sigma)}=-\sigma10^{-7}+\sigma2\pi j/16$.  Only the final
left/right regional charges are evaluated.  The same imposed record fixes the
total injected charge at every twist, so twist-dependent redistribution is
isolated by subtracting the corresponding infinitesimally displaced origin.
The experiment is a conditional static-response probe, not a quantized
monitored pump: neighboring twists are not dynamically connected and the
$2\pi$ endpoint must close by a large gauge transformation.  This note fixes
the scientific contract, offset rationale, charge estimator, numerical
checks, parallel execution, resumability, output schema, and claim boundary;
production plots and numerical conclusions are intentionally deferred.
\end{abstract}

\maketitle
\setcounter{tocdepth}{2}
\onecolumngrid
\tableofcontents
\twocolumngrid

\section{Purpose and status}

The adaptive preparation protocol of Bhuiyan, Pan, and Jian measures localized
overcomplete-Wannier (OW) occupations and applies outcome-dependent fermionic
feedback without postselection~\cite{BhuiyanPanJian2026}.  A twist around the
periodic $y$ direction probes transverse redistribution between two
$x$-separated interfaces, in the spirit of the Laughlin cylinder
argument~\cite{Laughlin1981,Halperin1982}.  Twisted boundary conditions also
provide a real-space formulation that does not require translation
symmetry~\cite{NiuThouless1984,NiuThoulessWu1985}.

This document describes a deliberately small, reproducible numerical pilot.
It records what is being calculated and how its correctness is established;
it does not report or interpret production curves.  The pilot is intended to
answer a narrow implementation question before a more expensive ensemble or
online-ramp calculation: if the stochastic word and therefore the net
feedback source are held fixed, does changing a static boundary twist alter
the final partition of charge between the two walls?

Four related constructions must not be conflated:

\begin{enumerate}
  \item The exact flattened-Hamiltonian benchmark continuously follows
  occupied wall branches and can exhibit nearly integer endpoint transfer.
  \item The present frozen-record experiment applies one imposed stochastic
  word in independent static simulations.
  \item An online monitored pump changes the twist during one uninterrupted
  Born-sampled trajectory and requires cycle-resolved source subtraction.
  \item Modular charge spreading evolves a localized packet in modular time;
  it is a propagation diagnostic rather than physical monitored time.
\end{enumerate}

Only item 2 is implemented here.

\section{Scientific contract}

\subsection{Geometry and OW measurements}

The physical system has $N_x\times N_y=16\times20$ unit cells with two
physical orbitals per cell.  A domain-wall mass profile places a topological
Chern-insulator region with $\alpha_1=1$ on $x\in[4,12]$ and a matched
trivial exterior with $\alpha_2=30$.  The OW basis uses trial family $X$ and
$n_{\mathrm{shell}}=1$.  Distinct OW stabilizers need not commute, and the OW
states form an overcomplete localized frame rather than an orthonormal Wannier
basis.

The two wall implementations are

\begin{itemize}
  \item \emph{soft}: \code{dw\_truncation=False} and
  \code{meas\_slab\_only=False}, retaining full OW tails and measurements
  throughout the system;
  \item \emph{hard}: \code{dw\_truncation=True} and
  \code{meas\_slab\_only=True}, terminating OW support and restricting
  measurements to the topological slab.
\end{itemize}

These are distinct stochastic protocols, so each receives its own reference
record.  They are compared at identical geometry, cycle count, initialization,
ordering, filling, and arithmetic precision.

\begin{table}[b]
\caption{Locked pilot configuration.}
\label{tab:contract}
\begin{ruledtabular}
\begin{tabular}{ll}
geometry & $N_x\times N_y=16\times20$ \\
domain wall & $x\in[4,12]$, $\alpha_1=1$, $\alpha_2=30$ \\
OW basis & trial $X$, $n_{\mathrm{shell}}=1$ \\
initial state & pure random, half filling \\
measurement order & random serial \\
feedback & perfect correction, no postselection \\
duration & $40=2N_y$ physical cycles \\
arithmetic & complex128 \\
twist points & $17$ per direction \\
independent records & one soft, one hard \\
canonical entry point & \code{classA\_U1FGTN.run\_markov\_circuit} \\
\end{tabular}
\end{ruledtabular}
\end{table}

\subsection{State representation and dynamics}

The canonical CPU engine stores the centered covariance $\Gamma$.  The
physical single-particle correlation matrix is
\begin{equation}
  G=\frac{\mathbf{1}+\Gamma}{2},
  \qquad
  G_{ij}=\operatorname{Tr}
  (\hat\rho\,\hat c_i^\dagger\hat c_j),
  \label{eq:correlation}
\end{equation}
where $\hat c_i$ is a physical lattice fermion and $\mathbf{1}$ is the
single-particle identity.  The runner calls
\code{classA\_U1FGTN.run\_markov\_circuit} for every sampled reference and
every forced replay.  It neither copies the Markov-cycle loop nor stores a
cycle-resolved covariance history.

Each trajectory begins from the same deterministic root seed and evolves for
40 complete cycles with pure-state rank-one covariance updates.  The reference
calculation records the ordered site identifiers, measurement outcomes,
conditional correction targets, and selected-branch probabilities required
for an exact replay.

\section{Reference record and static replay}

\subsection{Zero-twist record}

Let $\boldsymbol m_a$ denote the complete ordered word for arm
$a\in\{\mathrm{soft},\mathrm{hard}\}$.  It is generated by an ordinary
Born-sampled run at $\phi=0$.  The runner saves

\begin{equation}
  \bigl\{
  G_{a,\mathrm{init}}(0),
  G_{a,\mathrm{final}}(0),
  \boldsymbol m_a,
  \Delta N_{\mathrm{inj},a}
  \bigr\}.
\end{equation}

Before any twist tasks are admitted, a forced replay at $\phi=0$ must
reproduce the sampled final covariance.  The preregistered discrepancy is
\begin{equation}
  \epsilon_{0,a}
  =\frac{\lVert
  \Gamma_{a,\mathrm{replay}}(0)
  -\Gamma_{a,\mathrm{sampled}}(0)
  \rVert_F}{2N_xN_y}
  \le 10^{-12}.
  \label{eq:zero-parity}
\end{equation}
This check validates the record serialization, event ordering, initial state,
and deterministic replay path on a common input.

\subsection{Signed twist grid}

Write $\sigma=+1$ for increasing/counterclockwise flux and $\sigma=-1$ for
decreasing/clockwise flux.  The 17-point grid is
\begin{equation}
  \boxed{
  \phi_j^{(\sigma)}
  =-\sigma\varepsilon
   +\sigma\frac{2\pi j}{16}},
  \qquad
  j=0,\ldots,16,
  \qquad
  \varepsilon=10^{-7}.
  \label{eq:twist-grid}
\end{equation}
Thus the increasing sweep starts at $-10^{-7}$ and the decreasing sweep starts
at $+10^{-7}$.  The same exactly zero-twist reference word is used for both
directions; the offset changes the static evaluation point, not the sampled
record.

At each $\phi_j^{(\sigma)}$, the model and OW projectors are rebuilt at that
static twist.  The reference initial state is placed in the uniform gauge,
\begin{equation}
  G_{a,\mathrm{init}}(\phi)
  =U_\phi G_{a,\mathrm{init}}(0)U_\phi^\dagger,
  \qquad
  [U_\phi]_{xy\mu,xy\mu}
  =e^{\ii\phi y/N_y},
  \label{eq:gauge-init}
\end{equation}
and $\boldsymbol m_a$ is imposed through all 40 cycles.  Every twist is a
fresh static simulation.  There is no state transfer from point $j$ to
$j+1$.

\subsection{Why the \texorpdfstring{$10^{-7}$}{1e-7} offset is used}

The offset is not a physical mass, broadening, or Hamiltonian regulator.  It
resolves a convention at the exponentially small avoided crossing of the two
physical-wall branches at zero flux.  In the exact flattened-Hamiltonian
$20\times40$ benchmark, the closest levels at $\phi=0$ are approximately
\begin{align}
  E&=\pm5.49\times10^{-10} &&(\text{soft}),\\
  E&=\pm7.51\times10^{-10} &&(\text{hard}).
\end{align}
Starting infinitesimally to one side selects an unambiguous wall-localized
branch.  The quantization-favoring choice has sign opposite to the winding
direction, which is precisely Eq.~\eqref{eq:twist-grid}.  The exact continued
benchmark then gives

\begin{table}[b]
\caption{Continued exact-Hamiltonian reference at $20\times40$.  These values
motivate the signed origin convention; they are not acceptance targets for
the frozen-record circuit.}
\label{tab:exact}
\begin{ruledtabular}
\begin{tabular}{lrr}
wall & increasing from $-10^{-7}$ & decreasing from $+10^{-7}$ \\
\colrule
Soft & $+0.9999999976$ & $-0.9999999924$ \\
Hard & $+0.9999999987$ & $-0.9999999994$ \\
\end{tabular}
\end{ruledtabular}
\end{table}

Using the opposite side of the crossing gives magnitudes $0.98397184$ for the
soft wall and $0.98080969$ for the hard wall.  Those deficits reflect the
small initial weight on the opposite wall, not an inadequate flux grid.
Nearly integer transfer belongs to continuously tracked exact-Hamiltonian
branches.  Because the present pilot re-evaluates each twist independently,
quantization is explicitly excluded from its acceptance criteria.

\section{Endpoint charge observable}

\subsection{Fixed injected charge}

For a measurement event $e$, let $m_e\in\{0,1\}$ be the measured occupation
and $s_e\in\{0,1\}$ the perfect-correction target.  The imposed word fixes
the net particle source
\begin{equation}
  \Delta N_{\mathrm{inj},a}
  =\sum_{e\in\boldsymbol m_a}(s_e-m_e).
  \label{eq:injection}
\end{equation}
This integer is identical at every twist and for both orientations within the
same arm.  Each replay must satisfy global charge continuity,
\begin{equation}
  N_{\mathrm{final}}(\phi)
  -N_{\mathrm{initial}}(\phi)
  =\Delta N_{\mathrm{inj},a}
  \label{eq:global-continuity}
\end{equation}
to absolute tolerance $10^{-9}$.  The pilot deliberately measures charge only
after the entire record has been applied; no intermediate $\Delta N(\phi)$
is assigned.

\subsection{Left--right partition}

The cylinder is divided into two disjoint half-system regions,
\begin{align}
  R_L&:x=0,\ldots,7, & R_R&:x=8,\ldots,15,\\
  R_L+R_R&=\mathbf{1}.&&
  \label{eq:regions}
\end{align}
The endpoint regional charges are
\begin{equation}
  N_b(\phi)
  =\tr\!\left[
  R_bG_{a,\mathrm{final}}(\phi)
  \right],
  \qquad b=L,R.
\end{equation}
Each orientation is referenced to its own infinitesimally displaced origin,
\begin{align}
  \Delta N_b^{(\sigma)}(\phi_j)
  &=N_b(\phi_j^{(\sigma)})
    -N_b(\phi_0^{(\sigma)}),\\
  q_{\mathrm{wall}}^{(\sigma)}(\phi_j)
  &=\frac{
  \Delta N_R^{(\sigma)}(\phi_j)
  -\Delta N_L^{(\sigma)}(\phi_j)}{2}.
  \label{eq:qwall}
\end{align}
No regional source term is subtracted.  The global source is already identical
at all twists, so it cancels when the same-word origin is subtracted.  The
residual check
\begin{equation}
  \Delta N_L^{(\sigma)}
  +\Delta N_R^{(\sigma)}
  =0
  \label{eq:regional-balance}
\end{equation}
must hold to $10^{-9}$.

\subsection{Large-gauge closure}

The points $j=0$ and $j=16$ differ by a signed $2\pi$ flux and are related by
a large gauge transformation.  Endpoint covariances are retained only for
these two points, and the closure diagnostic is
\begin{equation}
  \epsilon_{\mathrm{cl}}^{(a,\sigma)}
  =\frac{
  \lVert
  U_{-\sigma2\pi}
  \Gamma_a(\phi_{16}^{(\sigma)})
  U_{-\sigma2\pi}^{\dagger}
  -\Gamma_a(\phi_0^{(\sigma)})
  \rVert_F}{2N_xN_y}.
  \label{eq:closure}
\end{equation}
The preregistered ceiling is $10^{-8}$.  A valid static replay should therefore
return $q_{\mathrm{wall}}(2\pi)$ to zero.  This closure is a correctness
condition and is the principal reason that the final point must not be read as
a quantized pump endpoint.

\section{Execution and resumability}

\subsection{Parallel task decomposition}

The full scan contains
\begin{equation}
  2\ \text{walls}
  \times2\ \text{directions}
  \times17\ \text{twists}
  =68
\end{equation}
deterministic replay tasks, plus two reference-record tasks.  The soft and hard
references are sampled in parallel.  Once both are verified, the 68 static
replays are distributed across a process pool.  The default host allocation
uses CPU cores 0--55, 56 workers, and one BLAS thread per worker.  Independent
twists of the same record are safe to parallelize because they share immutable
serialized inputs and no evolving state.

A single parent \code{tqdm} bar reports completed and failed twist points.
Worker stdout is suppressed to prevent 56 copies of the canonical engine
progress stream from making the log unreadable.  The campaign is launched
inside tmux so terminal detachment does not terminate the calculation.

\subsection{Completion-based resume}

Each twist task writes one compressed NPZ followed by one completion JSON.  The
completion file binds

\begin{itemize}
  \item the deterministic task and configuration hashes,
  \item the exact wall-specific record SHA-256,
  \item the result filename and byte count, and
  \item the result SHA-256.
\end{itemize}

On restart, a point is skipped only if the result and completion files both
exist and every declared identity, size, and checksum matches.  Missing,
partial, or mismatched pairs are rerun.  The expensive scientific unit is one
static replay, which is comfortably shorter than the repository threshold for
intra-task state checkpointing; therefore no cycle-level checkpoint, pointer
protocol, lease, Drive API, or distributed coordination layer is present.

The reference records use the same principle but bind three files: the
compressed trajectory word, the parent-state NPZ, and a completion JSON.  A
campaign-identity file prevents a changed configuration or source from being
silently resumed under the same identifier.

\section{Saved products and planned analysis}

The runner saves only what is required to reconstruct and audit the endpoint
response:

\begin{table}[b]
\caption{Pilot data products.}
\label{tab:products}
\begin{ruledtabular}
\begin{tabular}{lp{0.60\columnwidth}}
reference word & ordered measurement and correction events, gzip JSON \\
reference state & initial/final covariance, injection, replay diagnostics \\
twist point & $N_L,N_R,N_{\mathrm{tot}}$, $x$ profile, likelihood diagnostics \\
endpoint only & full covariance at $j=0,16$ for closure \\
completion & task/record/config identity, bytes, SHA-256 \\
aggregate & CSV and NPZ over all 68 points \\
manifest & configuration, sources, residual maxima, closure values \\
\end{tabular}
\end{ruledtabular}
\end{table}

Full covariance histories are deliberately excluded.  After all verified
points exist, the planned two-column diagnostic figure will show: (a) soft-wall
$\Delta N_L$ and $\Delta N_R$; (b) hard-wall $\Delta N_L$ and
$\Delta N_R$; (c) $q_{\mathrm{wall}}$ for both arms and directions; and
(d) $|\Delta N_L+\Delta N_R|$ against the preregistered tolerance.  No error
bars belong on the pilot curves because the 17 points are deterministic
functions of one shared word, not independent samples.

The first analysis pass must report, without threshold tuning,

\begin{enumerate}
  \item the two reference injections and zero-twist replay errors;
  \item the minimum selected-branch probability;
  \item maximum global-continuity and regional-balance residuals;
  \item four large-gauge closure errors;
  \item the peak and half-loop wall response for each arm and direction; and
  \item any disagreement between clockwise and counterclockwise static
  evaluations at gauge-equivalent twists.
\end{enumerate}

\section{Allowed interpretation}

If the structural checks pass and a finite $q_{\mathrm{wall}}(\phi)$ is
observed, the strongest justified pilot statement is:
\begin{quote}
For each of two explicit perfect-correction trajectories, the final Gaussian
state conditioned on the realized measurement word redistributes charge
between the two wall basins as a function of static boundary twist, while the
global feedback injection remains fixed.
\end{quote}

This is a trajectory-conditioned response statement.  It is not an ensemble
mean because there is only one independent record per wall.  It is not a
quantized pump because neighboring twists are independent static
recalculations and the $2\pi$ endpoint closes.  It is not evidence for a
physical current because the twist does not vary during monitored time.

A paper-level conditional-response result would repeat the full construction
for independent records and use the trajectory as the resampling unit.  A
physical monitored-pump claim would instead require one online Born trajectory
with a cycle-dependent twist, live charge accounting, and explicit
source-subtracted regional balance.  The present pilot should be used to
validate twist plumbing, record replay, wall partitioning, charge continuity,
and sign conventions before either extension.

\section{Reproducibility}

The self-contained project is
\path{00_WORKSPACE/CURRENT/frozen_record_flux_charge_pilot/}.
Its locked configuration is \code{campaign\_config.v1.json}; the campaign
runner is \code{run\_campaign.py}; and the local launcher is
\code{launch\_tmux.sh}.  The nominal campaign identifier is
\code{N16x20\_frozen\_flux\_charge\_v1}, with root seed 2026090202.
The result root is
\code{results/N16x20\_frozen\_flux\_charge\_v1/}.

The implementation has unit and small end-to-end CPU tests for the exact
68-task expansion, signed twist convention, injection parser, regional charge
partition, checksum failure, deterministic serial/parallel agreement, and
completion-based resume.  Production data are not part of this protocol note;
when incorporated later, their manifest and generated analysis inputs must be
treated as the numerical authority rather than values copied manually into
the manuscript.

\section{Conclusion}

This protocol isolates a useful but deliberately limited observable: the
endpoint charge response of a fixed monitored-circuit word to a static
boundary twist.  Its two wall constructions, signed $10^{-7}$ origins, fixed
injection identity, half-system partition, and large-gauge closure test are
fully specified before inspection of production curves.  The design is small,
parallel, and completion-resumable, while avoiding engineering layers that do
not protect scientific correctness.  Its eventual output will decide whether
the conditional response is promising enough to justify an ensemble
extension; it cannot by itself establish quantized or dynamical pumping.

\bibliographystyle{apsrev4-2}
\bibliography{references}

\end{document}
